\chapter{Apêndice C — Estrutura de Pastas e Reprodutibilidade}
\label{apendiceC}

Este apêndice descreve a organização dos arquivos e diretórios do repositório da
dissertação, bem como os princípios adotados para garantir reprodutibilidade
empírica, transparência metodológica e rastreabilidade dos resultados
apresentados nos capítulos principais.

\section{Estrutura de diretórios}

O repositório está organizado de forma modular, separando dados, scripts de
construção do \textit{pipeline}, scripts de regeneração de resultados, figuras,
tabelas e os próprios arquivos \LaTeX{} da dissertação. A estrutura abaixo
reproduz os principais diretórios e arquivos rastreados:

\begin{verbatim}
bitcoin-volatility-risk-premium/
│
├── data/                          # Conjuntos de dados
│   ├── vrp_with_targets.csv       # Dataset canônico (BVRP + retornos futuros)
│   ├── vrp_with_regimes.csv       # Dataset com classificação de regimes
│   ├── ml_dataset.csv             # Dataset dos modelos de ML (Cap. 6)
│   ├── btc_prices.csv             # Preços diários do Bitcoin (Binance)
│   ├── dvol_30d_full.csv          # IV30D — índice DVOL da Deribit
│   └── ...                        # Demais arquivos intermediários
│
├── code/                          # Pipeline de construção de dados
│   ├── download_dvol_deribit.py   # Download do DVOL via API Deribit
│   ├── update_btc_prices.py       # Atualização dos preços à vista
│   ├── build_vrp_dataset.py       # Cálculo de RV30D e construção do BVRP
│   ├── build_targets.py           # Retornos futuros e dataset canônico
│   ├── analyze_vrp_regimes.py     # Regime do BVRP (tercis, janela expansiva)
│   ├── build_ml_dataset.py        # Engenharia de features para ML
│   ├── analyze_magnitude_bvrp.py  # Teste de magnitude (apêndice do Cap. 8)
│   ├── analyze_bvrp_by_return_sign.py  # Análise por sinal do retorno
│   ├── run_all.py                 # Execução sequencial do pipeline
│   └── ...                        # Scripts de exploração e validação
│
├── scripts/                       # Regeneração de resultados por capítulo
│   ├── test_h1_bvrp_mean.py       # Teste de H1 com erro-padrão HAC (Cap. 5)
│   ├── regenerate_cap6_scatter.py # Figuras de previsão OOS (Cap. 6)
│   ├── regenerate_cap7.py         # Estratégias de negociação (Cap. 7)
│   ├── regenerate_cap7_figs.py    # Figuras complementares do Cap. 7
│   ├── regenerate_cap8.py         # Modelos supervisionados (Cap. 6)
│   ├── regenerate_cap9.py         # Análise de regimes (Cap. 8)
│   └── sync_overleaf.py           # Upload de figuras para o Overleaf
│
├── chapters/                      # Arquivos LaTeX dos capítulos
│   ├── cap1_introducao.tex
│   ├── cap2_referencial.tex
│   ├── cap3_dados.tex
│   ├── cap4_metodologia.tex
│   ├── cap5_resultados_ols.tex
│   ├── cap6_ml_unificado.tex
│   ├── cap7_estrategias.tex
│   ├── cap8_regimes.tex
│   ├── cap9_conclusao.tex
│   ├── Nota_metodologica.tex
│   ├── apendiceA_variaveis.tex
│   ├── apendiceB_scripts.tex
│   └── apendiceC_estrutura_pastas.tex
│
├── figs/                          # Figuras finais, por diretório de origem
│   ├── cap3/  cap5/  cap7/        # Descritivas, OLS e estratégias
│   ├── cap8/                      # Figuras do Cap. 6 (ML)
│   └── cap9/                      # Figuras do Cap. 8 (regimes)
│
├── tables/                        # Tabelas LaTeX finais
│   ├── cap3/  cap5/  tab7/        # Descritivas, OLS e estratégias
│   ├── tab8/                      # Tabelas do Cap. 6 (ML)
│   └── cap9/                      # Tabelas do Cap. 8 (regimes)
│
└── frontmatter/                   # Estrutura e preâmbulo do documento
    ├── main.tex                   # Arquivo raiz do documento LaTeX
    ├── preamble.tex               # Pacotes, margens e configurações
    ├── references.bib             # Referências bibliográficas (BibTeX)
    ├── resumo.tex   abstract.tex
    ├── capa.tex     folha_rosto.tex   folha_aprovacao.tex
    └── agradecimentos.tex   epigrafe.tex
\end{verbatim}

Cabe registrar que a nomenclatura de alguns diretórios e scripts precede a
reorganização dos capítulos: \texttt{regenerate\_cap8.py} gera os
resultados do Capítulo~\ref{chap:cap5_ml} e \texttt{regenerate\_cap9.py},
os do Capítulo~\ref{chap:regimes}, com os diretórios de saída
correspondentes (\texttt{tab8/}, \texttt{cap9/}). A correspondência é
indicada nos comentários acima.

\section{Fluxo de reprodutibilidade}

O \textit{pipeline} segue a seguinte ordem de execução:

\begin{enumerate}
    \item \textbf{Coleta de dados}: \texttt{code/download\_dvol\_deribit.py}
          obtém a série histórica do índice DVOL via API pública da Deribit e
          \texttt{code/update\_btc\_prices.py} atualiza os preços à vista via
          Binance.

    \item \textbf{Construção do BVRP}: \texttt{code/build\_vrp\_dataset.py}
          calcula a $RV_{30}$ a partir dos retornos logarítmicos diários,
          alinha com a $IV_{30}$ (DVOL) e gera o dataset principal
          \texttt{data/vrp\_30d\_dataset.csv}.

    \item \textbf{Variáveis-alvo}: \texttt{code/build\_targets.py} acrescenta
          os retornos futuros em seis horizontes
          ($h \in \{1, 5, 10, 20, 30, 60\}$ dias) e produz o dataset canônico
          \texttt{data/vrp\_with\_targets.csv}, utilizado nas análises dos
          Capítulos~\ref{chap:resultados_ols} e~\ref{chap:estrategias}.

    \item \textbf{Execução sequencial}: \texttt{code/run\_all.py} encadeia as
          etapas 1 a 3 --- coleta, construção do BVRP e geração das
          variáveis-alvo --- em uma única chamada. Os scripts de regime e de
          \textit{features} descritos a seguir são executados separadamente, na
          ordem indicada.

    \item \textbf{Regimes e \textit{features}}:
          \texttt{code/analyze\_vrp\_regimes.py} classifica o BVRP em tercis
          por janela expansiva, gerando \texttt{data/vrp\_with\_regimes.csv};
          \texttt{code/build\_ml\_dataset.py} constrói, a partir dele, as
          \textit{features} dos modelos supervisionados em
          \texttt{data/ml\_dataset.csv} (Capítulo~\ref{chap:cap5_ml}).

    \item \textbf{Regeneração por capítulo}: cada script em \texttt{scripts/}
          lê os datasets finais de \texttt{data/} e regenera todas as figuras
          (\texttt{figs/cap*/}) e tabelas (\texttt{tables/cap*/}) do capítulo
          correspondente, exportando os arquivos \texttt{.png} e \texttt{.tex}
          diretamente utilizados no documento.
\end{enumerate}

\section{Reprodutibilidade empírica}

Todos os valores numéricos reportados --- tabelas, coeficientes e métricas
--- são gerados por scripts determinísticos, com controle explícito de
sementes aleatórias (\texttt{random\_state=42}) nos modelos de aprendizado
de máquina. Figuras e tabelas são exportadas diretamente dos scripts Python
para os diretórios \texttt{figs/} e \texttt{tables/}, sem intervenções
manuais. Uma exceção conhecida diz respeito ao posicionamento gráfico dos
pontos no diagrama de enxame de valores SHAP, cujo deslocamento aleatório
não é fixado por semente: reexecuções produzem disposições visuais
ligeiramente distintas, sem alteração dos valores subjacentes.

A ordenação temporal é preservada em todas as etapas: a divisão
\textit{treino/teste} (70\,\%/30\,\%) respeita a cronologia das observações,
e os preditores são defasados quando necessário para assegurar ausência de
\textit{look-ahead bias}, conforme discutido nos
Capítulos~\ref{chap:metodologia} e~\ref{chap:cap5_ml}.

\section{Disponibilidade do código}

O repositório completo está disponível localmente e versionado com
\texttt{git}. Por razões de licenciamento de dados e uso de fontes externas
(API da Deribit e Binance), os dados brutos não são redistribuídos
publicamente; no entanto, os scripts de coleta permitem recriar a base a partir
das fontes originais, e todas as etapas metodológicas são descritas de forma
detalhada no texto principal e nos Apêndices A e B.

\section{Resumo}

Este apêndice documenta a estrutura computacional do projeto:
\begin{itemize}
    \item \textbf{\texttt{data/}} — datasets finais utilizados nas análises;
    \item \textbf{\texttt{code/}} — pipeline de coleta e construção das variáveis;
    \item \textbf{\texttt{scripts/}} — regeneração determinística de figuras e
          tabelas por capítulo;
    \item \textbf{\texttt{figs/}} e \textbf{\texttt{tables/}} — resultados
          organizados por capítulo e prontos para inclusão no documento;
    \item \textbf{\texttt{chapters/}} — arquivos \LaTeX{} dos capítulos,
          da nota metodológica e dos apêndices;
    \item \textbf{\texttt{frontmatter/}} — estrutura e preâmbulo do documento.
\end{itemize}

Essa organização reforça o compromisso com boas práticas de pesquisa empírica,
reprodutibilidade e transparência metodológica.
